% ============================================================
%  第二卷 第五章 恒定电磁场（§36–§45 + 习题）
%  底本：高教社中文版第8版；OCR 错误已修正
%  脚注已转录：5 处圈码已替换为 \footnote{...}（§37 两处、§41 三处；
%   转录自原书扫描页 p115/116/124/125，更新原第 16 条）
%  主要 OCR 修正：
%   1) \pmb{...}、{\cal E} 等粗体/花体混乱统一为 \boldsymbol；
%      花体 E 仅保留作能量 \mathcal{E}（§39 粒子能量、§52 场能量）
%   2) §37 “但 $\varphi$ 是我们知道”→“但是我们知道”（误植 φ）
%   3) §37 跨页断词“电动/力学”合并为“电动力学”
%   4) §38 习题解文末 OCR 乱码“LO”删去
%   5) §39 “如果 $M\boldsymbol{c}>|\alpha|$”粗体 c→正体 c
%   6) §39 “任意选择内 $\vec{\Gamma}$”→“任意选择内了”（OCR 乱码）
%   7) §39 习题2 偏转角 X→χ；“p 瑟福公式”→“卢瑟福公式”；衍字“的”删
%   8) §40 R_+ → R_{+-}（两电荷间径矢）
%   9) §41 (41.10) 加法定理中第二个极角补撇号 cos θ′（OCR 丢撇号）；
%      方位角统一记 Φ
%  10) §41 (41.15) 按定义 (41.11)(41.13) 端到端复核修正：
%      Q_±1^(2) = ∓(1/√6)(D_xz±iD_yz)（OCR 作 ±）；
%      Q_±2^(2) = −(1/(2√3))(D_xx−D_yy±2iD_xy)（OCR 作 1/(2√6)，
%      数值检验：单电荷置于 x 轴，Q_2 应为 −(√3/2)e）
%  11) §42 “式中，0是势在坐标原点之值”→φ_0；“电荷都 φ_0 集中”删误植；
%      “展开式（42.1）中的下一项”→（42.2）；2^l 极矩 D_m^(l)→Q_m^(l)；
%      系数 α_lm→a_lm（与 42.10 一致）
%  12) §43 矢量恒等式 OCR“cot f a”→“rot f a”
%  13) §44 “（略去指标 ^{a)}”→（略去指标 a）；Ā 式中乱码下标清除
%  14) §44 习题“我们得到 $\boldsymbol{p}$”末尾误植 p 删去
%  15) §45 (45.2) OCR“A·∇v”等→A·v＝(H×r)·v＝(r×v)·H
%  16) 脚注圈码①照 OCR 保留（§37 两处、§41 两处），脚注文本未录
%  17) 跨页断句标点补全若干，未改动文字
%  遗留问题：§37“Σe_a^+ = Σe_b^- = e”下标 b 照 OCR 保留；
%   §41 (41.10) 若原书确实排印两个 θ，则 θ′ 为按加法定理重构
% ============================================================
\chapter{恒定电磁场}

\section{库仑定律}

对于恒定电场，或者如通常所说静电场，麦克斯韦方程组有下面的形式：
\begin{equation}
  \operatorname{div}\boldsymbol{E}=4\pi\rho,
\end{equation}
\begin{equation}
  \operatorname{rot}\boldsymbol{E}=0.
\end{equation}
电场$\boldsymbol{E}$可以只利用一个标势来表示如下：
\begin{equation}
  \boldsymbol{E}=-\operatorname{grad}\varphi.
\end{equation}
将（36.3）代入（36.1），我们得到一个恒定电场的势所应满足的方程
\begin{equation}
  \Delta\varphi=-4\pi\rho.
\end{equation}
这个方程称为泊松方程.就特例言之，在真空中，即当 $\rho = 0$ 时，势满足拉普拉斯方程
\begin{equation}
  \Delta\varphi=0.
\end{equation}
从上面的方程可以得出一些结论，例如断定电场的势没有一处为最大或为最小.事实上，为要使 $\varphi$ 有极值，那么，$\varphi$ 对于坐标的一阶导数必须为零，而二阶导数 $\partial^2\varphi/\partial x^2$，$\partial^2\varphi/\partial y^2$，$\partial^2\varphi/\partial z^2$ 又都有同样的符号.后一要求是不可能的，因为满足了这样的要求，就不能满足（36.5）了.

我们现在来求一个点电荷所产生的场.从对称性的考虑可以知道，场的方向是从电荷 $e$ 的所在点沿着径矢指向空间各点.从同样的考虑还可以知道，场的大小 $E$ 仅与离开电荷的距离R有关.为了求出绝对值，我们应用方程(36.1)的积分形式(30.5).穿过一个半径为 R，包围着电荷e（以 $e$ 所在点为中心）的球面的电场通量等于 $4\pi R^2 E$；这个通量必须等于4πe.由此，我们得到
\begin{equation*}
  E=\frac{e}{R^2}.
\end{equation*}
用矢量符号表示，场$\boldsymbol{E}$可以写成
\begin{equation}
  \boldsymbol{E}=\frac{e\boldsymbol{R}}{R^3}.
\end{equation}
因此，一个点电荷所产生的场与从电荷算起的距离的平方成反比.就是库仑定律.这个场的势显然是
\begin{equation}
  \varphi=\frac{e}{R}.
\end{equation}
如果我们有一个电荷体系，那么，按照叠加原理，这个体系所产生的电场等于各个电荷单独所产生的电场之和.这个场的势是
\begin{equation*}
  \varphi=\sum_a\frac{e_a}{R_a},
\end{equation*}
式中 $R_a$ 是从电荷 $e_a$ 的所在点到我们求它的势的那一点的距离.如果我们引用电荷密度 $\rho$，这个公式就变为
\begin{equation}
  \varphi=\int\frac{\rho}{R}\,\mathrm{d}V,
\end{equation}
其中 R 是从体积元 $\mathrm{d}V$ 到给定点的距离.

将点电荷的 $\rho$ 及 $\varphi$ 的值，即 $\rho = e\delta(R)$ 及 $\varphi = e/R$，代入（36.4）式，我们得到
\begin{equation}
  \Delta\frac{1}{R}=-4\pi\delta(R).
\end{equation}

\section{电荷的静电能}

我们来求一个电荷体系的能量.我们将从场的能量的概念出发，即从能量密度的公式(31.5）出发来讨论.一个电荷体系的能量应当等于
\begin{equation*}
  U=\frac{1}{8\pi}\int E^2\,\mathrm{d}V,
\end{equation*}
其中E是诸电荷所产生的场，而积分则应该遍及全部空间.将 $\boldsymbol{E}=-\operatorname{grad}\varphi$ 代入，则可按下面的方式来变换U：
\begin{equation*}
  U=-\frac{1}{8\pi}\int\boldsymbol{E}\cdot\operatorname{grad}\varphi\,\mathrm{d}V
  =-\frac{1}{8\pi}\int\operatorname{div}(\boldsymbol{E}\varphi)\,\mathrm{d}V
  +\frac{1}{8\pi}\int\varphi\operatorname{div}\boldsymbol{E}\,\mathrm{d}V.
\end{equation*}
按照高斯定理，第一个积分等于 $\boldsymbol{E}\varphi$ 在包围积分体积的曲面上的积分，但是因为场在无穷远处为零，而积分又要遍及全部空间，所以第一个积分为零.将 $\operatorname{div}\boldsymbol{E}=4\pi\rho$ 代入第二个积分，则得到电荷体系的能量表达式
\begin{equation}
  U=\frac{1}{2}\int\rho\varphi\,\mathrm{d}V.
\end{equation}
对于一个点电荷体系 $e_a$，我们可以用对电荷求和的符号来代替积分号：
\begin{equation}
  U=\frac{1}{2}\sum e_a\varphi_a,
\end{equation}
其中，$\varphi_a$ 是所有电荷在 $e_a$ 所在点所产生的场的势.

按照库仑定律，如果我们对一个带电的基本粒子（例如电子）和粒子本身所产生的场应用所得的公式，那么，我们就得到一个结论，即电荷应当具有等于 $e\varphi/2$ 的“自己的”势能，这里$\varphi$是电荷在其所在点产生的场的势.但是我们知道，在相对论中，每一个基本粒子应当当做一个点.因此，基本粒子的场的势 $\varphi = e/R$ 在 $R = 0$ 的这一点变为无穷大.这样一来，按照电动力学，电子就应当具有无限大的“自”能，因而也就有无限大的质量.这个结论的物理荒谬性表明，电动力学本身的基本原理就导致一个结果，即电动力学的应用应当限制在一定范围内.

我们要注意，从电动力学，我们得到了无限大的“自”能与质量，因此在电动力学中就不能提出到底电子的总质量是不是电磁质量（就是与粒子的电磁自能有关的质量）这个问题.\footnote{从纯形式的观点看，电子质量的有限性可以通过如下方式来处理，即引入一个非电磁起源的无限大负质量来补偿电磁质量的无限性（质量“重正化”）.然而，我们下面将会看到（§75），这并不能消除经典电动力学的所有内在矛盾.}

既然基本粒子的没有物理意义的无限大“自”能的出现，与粒子应当看做一个点的事实有关，那么我们可以得出结论：作为一个逻辑上完备的物理理论的电动力学，当过渡到充分小的距离时，就成为自相矛盾的了.我们可以提出这个距离有多大数量级的问题.注意到电子的电磁自能与静能 $mc^2$ 同数量级，就可以答复这个问题了.另外一方面，如果我们认为电子具有一定的半径 $R_0$，那么，它的自有势能应该与 $e^2/R_0$ 同数量级.从这两个量应同级的要求，亦即从 $e^2/R_0\sim mc^2$ 的要求出发，我们得到
\begin{equation}
  R_0\sim\frac{e^2}{mc^2}.
\end{equation}
这个尺度（称为电子的“半径”）决定了电动力学适用于电子的范围，这是从它的基本原理得来的.但是我们应该注意，实际上，由于量子现象，电动力学应用范围比我们在这里所确定的范围还要小得多\footnote{量子效应在与 $h/(mc)$ 同数量级的距离上就变得很重要了，这里的 $h$ 是普朗克常数.这些距离同 $R_0$ 之比量级为 $hc/e^2\sim 137$.}.

现在我们再回到公式(37.2).从库仑定律知道，公式中的势 $\varphi_a$ 等于
\begin{equation}
  \varphi_a=\sum\frac{e_b}{R_{ab}},
\end{equation}
其中，$R_{ab}$ 是电荷 $e_a,e_b$ 间的距离.能量的表达式(37.2)包含两部分.第一，它包含一个无限大的常数，即电荷的自能，它与电荷之间的相互位置无关.第二部分是电荷的相互作用能，它与电荷之间的相互位置有关.显然，物理学上关心的仅仅是第二部分.这部分等于
\begin{equation}
  U'=\frac{1}{2}\sum e_a\varphi_a',
\end{equation}
其中
\begin{equation}
  \varphi_a'=\sum_{b(\neq a)}\frac{e_b}{R_{ab}}
\end{equation}
是除了 $e_a$ 以外的所有电荷在 $e_a$ 所在点产生的势.换句话说，我们可以写
\begin{equation}
  U'=\frac{1}{2}\sum_{a\neq b}\frac{e_ae_b}{R_{ab}}.
\end{equation}
就特例言之，两个粒子的相互作用能是
\begin{equation}
  U'=\frac{e_1e_2}{R_{12}}.
\end{equation}

\section{匀速运动电荷的场}

我们来求一个以速度V匀速运动的电荷e所产生的场.我们称静止参考系为K系；称随电荷一起运动的参考系为 $K'$ 系.设电荷位于 $K'$ 系的坐标原点，$K'$ 系沿x轴对K作相对运动，y轴与z轴分别平行于 $y'$ 与 $z'$.在 $t = 0$ 的时刻，两个系统的原点相重合.因此，电荷在K系中的坐标是 $x = Vt,\ y = z = 0$.在 $K'$ 系中，我们有一个恒定电场，它的矢势 $\boldsymbol{A}'=\boldsymbol{0}$，而它的标势则等于 $\varphi' = e/R'$，式中 $R'^2 = x'^2 + y'^2 + z'^2$.按照（24.1），当 $\boldsymbol{A}'=0$ 时，在K系中，
\begin{equation}
  \varphi=\frac{\varphi'}{\sqrt{1-\dfrac{V^2}{c^2}}}
  =\frac{e}{R'\sqrt{1-\dfrac{V^2}{c^2}}}.
\end{equation}
现在我们应当用K系中的 $x,y,z$ 来表示 $R'$.按照洛伦兹变换公式，
\begin{equation*}
  x'=\frac{x-Vt}{\sqrt{1-\dfrac{V^2}{c^2}}},\qquad
  y'=y,\qquad z'=z,
\end{equation*}
由此得到
\begin{equation}
  R'^2=\frac{(x-Vt)^2+\left(1-\dfrac{V^2}{c^2}\right)(y^2+z^2)}{1-\dfrac{V^2}{c^2}}.
\end{equation}
将此式代入(38.1)，就得到
\begin{equation}
  \varphi=\frac{e}{R^*},
\end{equation}
式中，我们引入了记号 $R^*$
\begin{equation}
  R^{*2}=(x-Vt)^2+\left(1-\frac{V^2}{c^2}\right)(y^2+z^2).
\end{equation}
在K系中，矢势等于
\begin{equation}
  \boldsymbol{A}=\varphi\frac{\boldsymbol{V}}{c}=\frac{e\boldsymbol{V}}{cR^*}.
\end{equation}
在 $K'$ 系中，磁场 $\boldsymbol{H}'$ 不存在，而电场则是
\begin{equation*}
  \boldsymbol{E}'=\frac{e\boldsymbol{R}'}{R'^3}.
\end{equation*}
从公式（24.2），我们得到
\begin{equation*}
\begin{aligned}
  &E_x=E_x'=\frac{ex'}{R'^3},\qquad
  E_y=\frac{E_y'}{\sqrt{1-\dfrac{V^2}{c^2}}}
  =\frac{ey'}{R'^3\sqrt{1-\dfrac{V^2}{c^2}}},\\
  &E_z=\frac{ez'}{R'^3\sqrt{1-\dfrac{V^2}{c^2}}},
\end{aligned}
\end{equation*}
将用 $x,y,z$ 表示的 $R',x',y',z'$ 的式子代入，就得到
\begin{equation}
  \boldsymbol{E}=\left(1-\frac{V^2}{c^2}\right)\frac{e\boldsymbol{R}}{R^{*3}},
\end{equation}
式中，R 是从电荷 e到坐标为 $x,y,z$ 的点的径矢（它的分量是 $x-Vt,\ y,\ z$）.

如果我们引入运动方向与径矢R所夹之角θ,E的表达式可以写成另一形式.显然，$y^2+z^2=R^2\sin^2\theta$，所以 $R^{*2}$ 可以写成下面的形式：
\begin{equation}
  R^{*2}=R^2\left(1-\frac{V^2}{c^2}\sin^2\theta\right).
\end{equation}
于是，对于 $\boldsymbol{E}$，我们就有
\begin{equation}
  E=\frac{eR}{R^3}\,
  \frac{1-\dfrac{V^2}{c^2}}{\left(1-\dfrac{V^2}{c^2}\sin^2\theta\right)^{3/2}}.
\end{equation}
在离电荷的距离为R时，若θ从零增加到 $\pi/2$（或者θ从π减至 $\pi/2$）则场E的值将随之增加.沿着运动方向的场 $E_{\parallel}$（θ=0 或π）将有最小值；它等于
\begin{equation*}
  E_{\parallel}=\frac{e}{R^2}\left(1-\frac{V^2}{c^2}\right).
\end{equation*}
最大的场与速度垂直（$\theta=\pi/2$），并等于
\begin{equation*}
  E_{\perp}=\frac{e}{R^2}\frac{1}{\sqrt{1-\dfrac{V^2}{c^2}}}.
\end{equation*}
我们要注意，当速度增加时，场 $E_{\parallel}$ 减小，而 $E_{\perp}$ 则增大.我们可以形象地叙述这种现象：一个运动电荷的电场在运动的方向“收缩”.对于一个同光速接近的速度V，公式（38.8)中的分母在 $\theta=\pi/2$ 的附近的一个狭小范围内接近于零，这个范围的“宽度”的数量级为
\begin{equation*}
  \Delta\theta\sim\sqrt{1-\frac{V^2}{c^2}}.
\end{equation*}
因此，仅仅在赤道平面的邻近的一个狭小角度范围内，一个快速运动着的电荷的电场是大的，而这个范围则随着V的增加而减小，其减小的情况如同 $\sqrt{1-V^2/c^2}$.

在K系内，磁场等于
\begin{equation}
  \boldsymbol{H}=\frac{1}{c}\boldsymbol{V}\times\boldsymbol{E}
\end{equation}
（见（24.5））.在 $V\ll c$ 的情况下，我们近似地得到 $\boldsymbol{E}=e\boldsymbol{R}/R^3$，而磁场为
\begin{equation}
  \boldsymbol{H}=\frac{e}{c}\frac{\boldsymbol{V}\times\boldsymbol{R}}{R^3}.
\end{equation}

\begin{xiti}

\noindent{\heiti 习\ 题}\quad 求两个以相同速度V运动的电荷之间的力（在K系中）.

解：我们要求的力F可以通过计算一个电荷（$e_1$）在另一个电荷（$e_2$）产生的场中所受到的力来决定.用(38.9)式，我们得到
\begin{equation*}
  \boldsymbol{F}=e_1\boldsymbol{E}_2+\frac{e_1}{c}\boldsymbol{V}\times\boldsymbol{H}_2
  =e_1\left(1-\frac{V^2}{c^2}\right)\boldsymbol{E}_2
  +\frac{e_1}{c^2}\boldsymbol{V}(\boldsymbol{V}\cdot\boldsymbol{E}_2).
\end{equation*}
由（38.8）代入 $\boldsymbol{E}_2$，我们就得到力在运动方向的分量 $(F_x)$ 和垂直于它的分量 $(F_y)$：
\begin{equation*}
  F_x=\frac{e_1e_2}{R^2}
  \frac{\left(1-\dfrac{V^2}{c^2}\right)\cos\theta}{\left(1-\dfrac{V^2}{c^2}\sin^2\theta\right)^{3/2}},\qquad
  F_y=\frac{e_1e_2}{R^2}
  \frac{\left(1-\dfrac{V^2}{c^2}\right)^2\sin\theta}{\left(1-\dfrac{V^2}{c^2}\sin^2\theta\right)^{3/2}},
\end{equation*}
式中R是从 $e_2$ 到 $e_1$ 的径矢，θ是R和V之间的夹角.

\end{xiti}

\section{库仑场内的运动}

我们来研究一个质量为 $m$、电荷为e的粒子在另一个电荷 $e'$ 所产生的场内的运动；我们假设第二个电荷的质量比 $m$ 大得如此之多，以致我们可以把它当做固定的.这样一来，我们的问题就化成了研究一个电荷 $e$ 在中心对称电场（其势为 $\varphi = e'/r$）内的运动.

粒子的总能量 $\mathcal{E}$ 等于
\begin{equation*}
  \mathcal{E}=c\sqrt{p^2+m^2c^2}+\frac{\alpha}{r},
\end{equation*}
式中，$\alpha = ee'$.如果我们在粒子的运动平面内采用极坐标，那么，由力学可知
\begin{equation*}
  p^2=\frac{M^2}{r^2}+p_r^2,
\end{equation*}
式中的 $p_r$ 是动量的径向分量，而M则是粒子的恒定角动量.这时，
\begin{equation}
  \mathcal{E}=c\sqrt{p_r^2+\frac{M^2}{r^2}+m^2c^2}+\frac{\alpha}{r}.
\end{equation}
现在我们来讨论粒子在运动过程中能否随意地接近中心的问题.首先，很容易看出，如果 $e$ 同 $e'$ 相互排斥，即 $e$ 同 $e'$ 有同样的符号，那么，这种接近就永远不可能.此外，在相互吸引的情形下(e同 $e'$ 有相反的符号)，如果 $Mc>|\alpha|$，随意接近中心是不可能的，因为在这种情况下，(39.1)式中的第一项永比第二项大，并且当 $r\to 0$ 时，方程的右边就要变为无限大.反之，如果 $Mc<|\alpha|$，那么，当 $r\to 0$ 时，这个式子可以保持有限值（在此，不用说，$p_r$ 趋向无限大).因此，如果
\begin{equation}
  Mc<|\alpha|,
\end{equation}
粒子在运动过程中就能够“降落”到吸引它的电荷上，这与非相对论力学不同，在非相对论力学中，对于库仑场，这样的“降落”一般是不可能的（只有$M = 0$ 情形除外，这时粒子 $e$ 沿着一条直线飞向 $e'$）.

为了完全决定一个电荷在库仑场内的运动，从哈密顿-雅可比方程出发是最便利的.我们在运动平面内取极坐标 $r,\varphi$.哈密顿-雅可比方程（16.11）可以写成
\begin{equation*}
  -\frac{1}{c^2}\left(\frac{\partial S}{\partial t}+\frac{\alpha}{r}\right)^2
  +\left(\frac{\partial S}{\partial r}\right)^2
  +\frac{1}{r^2}\left(\frac{\partial S}{\partial\varphi}\right)^2
  +m^2c^2=0.
\end{equation*}
我们来寻求下面形式的S:
\begin{equation*}
  S=-\mathcal{E}t+M\varphi+f(r),
\end{equation*}
式中，$\mathcal{E}$ 及M分别为运动粒子的恒定能量及恒定角动量.结果我们求得
\begin{equation}
  S=-\mathcal{E}t+M\varphi
  +\int\sqrt{\frac{1}{c^2}\left(\mathcal{E}-\frac{\alpha}{r}\right)^2
  -\frac{M^2}{r^2}-m^2c^2}\,\mathrm{d}r.
\end{equation}
轨道由方程 $\partial S/\partial M=\mathrm{const}$ 决定.积分（39.3）后，得到下面的结果：

(a) 如果 $Mc>|\alpha|$
\begin{equation}
\begin{aligned}
  \left(c^2M^2-\alpha^2\right)\frac{1}{r}={}&
  c\sqrt{(M\mathcal{E})^2-m^2c^2(M^2c^2-\alpha^2)}
  \cos\left(\varphi\sqrt{1-\frac{\alpha^2}{c^2M^2}}\right)\\
  &-\mathcal{E}\alpha;
\end{aligned}
\end{equation}

(b) 如果 $Mc<|\alpha|$
\begin{equation}
\begin{aligned}
  \left(\alpha^2-c^2M^2\right)\frac{1}{r}={}&
  \pm c\sqrt{(M\mathcal{E})^2+m^2c^2(\alpha^2-M^2c^2)}
  \cosh\left(\varphi\sqrt{\frac{\alpha^2}{c^2M^2}-1}\right)\\
  &+\mathcal{E}\alpha;
\end{aligned}
\end{equation}

(c) 如果 $Mc=|\alpha|$
\begin{equation}
  \frac{2\mathcal{E}\alpha}{r}
  =\mathcal{E}^2-m^2c^4-\varphi^2\left(\frac{\mathcal{E}\alpha}{cM}\right)^2.
\end{equation}
积分常数已经包含在角 $\varphi$ 的计算起点的任意选择内了.

在(39.4)式中，平方根前面未指明正负号并不重要，因为余弦符号后的角$\varphi$ 的计算起点的选择是任意的.在相互吸引的情形下（$\alpha < 0$），如果 $\mathcal{E}<mc^2$ 与这个方程相对应的轨道上，r之值皆为有限(有限运动).如果 $\mathcal{E}>mc^2$，那么，r可以是无限大（无限运动).在非相对论力学中，有限运动与在闭合轨道（椭圆）上的运动相对应.从（39.4）式可以看出，在相对论力学中，轨道不可能是封闭的；当 $\varphi$ 变动 $2\pi$ 时，到中心的距离 $r$ 并不回到它的原来的值.我们得到的轨道不是椭圆，而是张开的“玫瑰形”，因此，在非相对论力学中，在库仑场中的有限运动的轨道是封闭的，而在相对论力学中，库仑场就失去了这个特性.

在（39.5）式中，当 $\alpha < 0$ 时，根号前应当选择正号，而当 $\alpha > 0$ 时，就应当选择负号（正负号的另一种选择对应于(39.1）式中的根号前正负号改变).

在 $\alpha < 0$ 的情况下，轨道（39.5）及（39.6）是螺线，当 $\varphi\to\infty$ 时，距离 $r\to 0$.电荷“降落”到坐标原点所需要的时间是有限的.这可以从下面看出，r的坐标与时间的关系是由方程 $\partial S/\partial\mathcal{E}=\mathrm{const}$ 决定的；将（39.3）代入，我们看出，决定时间的积分当 $r\to 0$ 时是收敛的.

\begin{xiti}

\noindent{\heiti 习题1}\quad 求一个电荷在排斥的库仑场（$\alpha > 0$）中飞行时的偏转角.

解：偏转角 $\chi = \pi - 2\varphi_0$，这里 $2\varphi_0$ 是轨道（39.4）的两个渐近线所夹之角.我们求得
\begin{equation*}
  \chi=\pi-\frac{2cM}{\sqrt{c^2M^2-\alpha^2}}
  \arctan\frac{v\sqrt{c^2M^2-\alpha^2}}{c\alpha},
\end{equation*}
式中，v是电荷在无穷远处的速度.

\noindent{\heiti 习题2}\quad 求质点在库仑场中小偏角散射的有效截面.

解：有效截面dσ是在一秒钟内散射到一定的立体角元 $\mathrm{d}o$ 内的粒子数与被散射粒子的通量密度(即每秒经过垂直于粒子束的 $1\,\mathrm{cm}^2$ 面积的粒子数）之比.

因为粒子在通过场时的偏转角χ由碰撞参量 $\rho$ 决定（$\rho$ 即从中心到一条直线的距离，这条直线就是粒子在场不存在时应该沿其运动的直线)，
\begin{equation*}
  \mathrm{d}\sigma=2\pi\rho\,\mathrm{d}\rho
  =2\pi\rho\frac{\mathrm{d}\rho}{\mathrm{d}\chi}\,\mathrm{d}\chi
  =\rho\frac{\mathrm{d}\rho}{\mathrm{d}\chi}\frac{\mathrm{d}o}{\sin\chi},
\end{equation*}
式中，$\mathrm{d}o=2\pi\sin\chi\,\mathrm{d}\chi$ (参看本教程第一卷§18).偏转角(如果很小)可以认为等于动量的变化对于它的初值之比.动量的变化等于作用于电荷上的力的时间积分，这个力在垂直于运动方向的分量近似地等于 $(\alpha/r^2)\cdot(\rho/r)$.因此，我们得到
\begin{equation*}
  \chi=\frac{1}{p}\int_{-\infty}^{+\infty}
  \frac{\alpha\rho\,\mathrm{d}t}{(\rho^2+v^2t^2)^{3/2}}
  =\frac{2\alpha}{p\rho v}
\end{equation*}
(v 是粒子的速度).从而，我们求得在小 $\chi$ 情况下的有效截面：
\begin{equation*}
  \mathrm{d}\sigma=4\left(\frac{\alpha}{pv}\right)^2\frac{\mathrm{d}o}{\chi^4}.
\end{equation*}
在非相对论情形下，$p\approx mv$，这个表达式与从卢瑟福公式在小 $\chi$ 时得到的一致（参见本教程第一卷§19).

\end{xiti}

\section{偶极矩}

我们现在来研究一个电荷体系在与这个体系相距很远之处所产生的场，所谓很远之处，就是说该处与体系的距离比这个体系中各电荷间的距离大得多.

我们引入一个坐标系，它的原点在电荷体系内的任意一点上.设各电荷的径矢为 $\boldsymbol{r}_a$.所有电荷在以 $R_0$ 为径矢的点所生的场的势等于
\begin{equation}
  \varphi=\sum\frac{e_a}{|\boldsymbol{R}_0-\boldsymbol{r}_a|}
\end{equation}
(求和遍及所有电荷)；式中，$\boldsymbol{R}_0-\boldsymbol{r}_a$ 是从电荷 $e_a$ 到我们正在求势那一点的径矢.

我们必须对大的 $R_0\ (R_0\gg r_a)$ 来研究这个表达式.为此，我们将它展开为 $r_a/R_0$ 的幂级数，利用公式
\begin{equation*}
  f(\boldsymbol{R}_0-\boldsymbol{r})=f(\boldsymbol{R}_0)
  -\boldsymbol{r}\cdot\operatorname{grad}f(\boldsymbol{R}_0)
\end{equation*}
(梯度符号 grad 是对矢量 $R_0$ 的端点坐标进行微分）.准确到一级项，
\begin{equation}
  \varphi=\frac{\sum e_a}{R_0}
  -\sum e_a\boldsymbol{r}_a\cdot\operatorname{grad}\frac{1}{R_0}.
\end{equation}
我们称
\begin{equation}
  \boldsymbol{d}=\sum e_a\boldsymbol{r}_a
\end{equation}
为电荷体系的偶极矩，应该注意，如果所有的电荷之和为零，即 $\sum e_a = 0$ 那么，偶极矩就与坐标原点的选择无关，因为同一个电荷在两个不同的坐标系中的径矢 $\boldsymbol{r}_a$ 及 $\boldsymbol{r}_a'$ 有下面的关系：
\begin{equation*}
  \boldsymbol{r}_a'=\boldsymbol{r}_a+\boldsymbol{a},
\end{equation*}
式中，$\boldsymbol{a}$ 是一个常矢量.因此，如果 $\sum e_a = 0$，偶极矩在两个坐标系中是相等的：
\begin{equation*}
  \boldsymbol{d}'=\sum e_a\boldsymbol{r}_a'
  =\sum e_a\boldsymbol{r}_a+\boldsymbol{a}\sum e_a=\boldsymbol{d}.
\end{equation*}
如果我们用 $e_a^+,\boldsymbol{r}_a^+$ 和 $-e_a^-,\boldsymbol{r}_a^-$ 代表这个体系的正电荷和负电荷以及它们的径矢，那么，我们可以将偶极矩写成
\begin{equation}
  \boldsymbol{d}=\sum e_a^+\boldsymbol{r}_a^+-\sum e_a^-\boldsymbol{r}_a^-
  =\boldsymbol{R}^+\sum e_a^+-\boldsymbol{R}^-\sum e_a^-,
\end{equation}
式中，
\begin{equation}
  \boldsymbol{R}^+=\frac{\sum e_a^+\boldsymbol{r}_a^+}{\sum e_a^+},\qquad
  \boldsymbol{R}^-=\frac{\sum e_a^-\boldsymbol{r}_a^-}{\sum e_a^-}
\end{equation}
是正电荷及负电荷的“电荷中心”的径矢.如果 $\sum e_a^+ = \sum e_b^- = e$，那么，
\begin{equation}
  \boldsymbol{d}=e\boldsymbol{R}_{+-},
\end{equation}
式中，$\boldsymbol{R}_{+-}=\boldsymbol{R}^+-\boldsymbol{R}^-$ 是从正电荷中心到负电荷中心的径矢.就特殊情形言之，如果只有两个电荷，那么，$\boldsymbol{R}_{+-}$ 就是这两个电荷间的径矢.

如果 $\sum e_a = 0$，那么，这个体系在远距离处的场的势是：
\begin{equation}
  \varphi=-\boldsymbol{d}\cdot\nabla\frac{1}{R_0}
  =\frac{\boldsymbol{d}\cdot\boldsymbol{R}_0}{R_0^3}.
\end{equation}
场强E是：
\begin{equation*}
  \boldsymbol{E}=-\operatorname{grad}\frac{\boldsymbol{d}\cdot\boldsymbol{R}_0}{R_0^3}
  =-\frac{1}{R_0^3}\operatorname{grad}(\boldsymbol{d}\cdot\boldsymbol{R}_0)
  -(\boldsymbol{d}\cdot\boldsymbol{R}_0)\operatorname{grad}\frac{1}{R_0^3},
\end{equation*}
或者，最后有
\begin{equation}
  \boldsymbol{E}=\frac{3(\boldsymbol{n}\cdot\boldsymbol{d})\boldsymbol{n}-\boldsymbol{d}}{R_0^3},
\end{equation}
式中 $\boldsymbol{n}$ 是沿 $\boldsymbol{R}_0$ 的单位矢量.场的另一个有用表达式为
\begin{equation}
  \boldsymbol{E}=(\boldsymbol{d}\cdot\nabla)\nabla\frac{1}{R_0}.
\end{equation}
因此，一个总电荷等于零的电荷体系在远距离处所产生的场的势与到这个体系的距离的平方成反比，而场的强度则与距离的立方成反比.这个场围绕d的方向具有轴对称性.在穿过该方向（我们取为z轴）的平面内，矢量E的分量是：
\begin{equation}
  E_z=d\frac{3\cos^2\theta-1}{R_0^3},\qquad
  E_x=d\frac{3\sin\theta\cos\theta}{R_0^3}.
\end{equation}
在这个平面内的径向和切向分量是
\begin{equation}
  E_R=d\frac{2\cos\theta}{R_0^3},\qquad
  E_\theta=-d\frac{\sin\theta}{R_0^3}.
\end{equation}

\section{多极矩}

在势按 $1/R_0$ 的幂展开的展开式
\begin{equation}
  \varphi=\varphi^{(0)}+\varphi^{(1)}+\varphi^{(2)}+\cdots
\end{equation}
中，$\varphi^{(n)}$ 这一项与 $1/R_0^{n+1}$ 成正比.我们看出，第一项 $\varphi^{(0)}$ 为所有电荷的和所决定；第二项 $\varphi^{(1)}$ 有时称为体系的偶极势，为体系的偶极矩所决定.

展开式中的第三项是
\begin{equation}
  \varphi^{(2)}=\frac{1}{2}\sum ex_\alpha x_\beta
  \frac{\partial^2}{\partial X_\alpha\partial X_\beta}\frac{1}{R_0},
\end{equation}
式中的求和遍历所有电荷，我们在此没有写出电荷的编号；$x_\alpha$ 是矢量r的分量，$X_\alpha$ 是矢量 $\boldsymbol{R}_0$ 的分量.势的这一部分通常称为四极势.如果体系的电荷的和及偶极矩的和都等于零，展开式就从 $\varphi^{(2)}$ 开始.

在（41.2）式中，有六个量 $\sum ex_\alpha x_\beta$ 出现.但是，很容易看出场实际上并不与六个独立的量有关，而仅与五个独立的量有关.这是因为函数 $1/R_0$ 满足拉普拉斯方程，即
\begin{equation*}
  \Delta\frac{1}{R_0}\equiv\delta_{\alpha\beta}
  \frac{\partial^2}{\partial X_\alpha\partial X_\beta}\frac{1}{R_0}=0.
\end{equation*}
因而，我们可以将 $\varphi^{(2)}$ 写成下面的形式：
\begin{equation*}
  \varphi^{(2)}=\frac{1}{2}\sum e\left(x_\alpha x_\beta-\frac{1}{3}r^2\delta_{\alpha\beta}\right)
  \frac{\partial^2}{\partial X_\alpha\partial X_\beta}\frac{1}{R_0}.
\end{equation*}
张量
\begin{equation}
  D_{\alpha\beta}=\sum e(3x_\alpha x_\beta-r^2\delta_{\alpha\beta})
\end{equation}
称为这个体系的四极矩.从 $D_{\alpha\beta}$ 的定义，很容易看出，对角线上的元素之和为零：
\begin{equation}
  D_{\alpha\alpha}=0.
\end{equation}
因此对称张量 $D_{\alpha\beta}$ 一共有五个独立的分量.利用 $D_{\alpha\beta}$，我们可以写出
\begin{equation}
  \varphi^{(2)}=\frac{D_{\alpha\beta}}{6}
  \frac{\partial^2}{\partial X_\alpha\partial X_\beta}\frac{1}{R_0},
\end{equation}
或者，进行微分，
\begin{equation*}
  \frac{\partial^2}{\partial X_\alpha\partial X_\beta}\frac{1}{R_0}
  =\frac{3X_\alpha X_\beta}{R_0^5}-\frac{\delta_{\alpha\beta}}{R_0^3},
\end{equation*}
并考虑到 $\delta_{\alpha\beta}D_{\alpha\beta}=D_{\alpha\alpha}=0$
\begin{equation}
  \varphi^{(2)}=\frac{D_{\alpha\beta}n_\alpha n_\beta}{2R_0^3}.
\end{equation}
像所有对称的三维张量一样，张量 $D_{\alpha\beta}$ 可以化到主轴.由于(41.4)，三个主值中一般只有两个是独立的.如果正好电荷系统围绕某个轴（z轴）对称\footnote{指的是任何高于 2 阶的对称轴.}，则这个轴必为张量 $D_{\alpha\beta}$ 的主轴之一，其他两个轴在 $xy$ 平面内的位置是任意的，三个主值之间的关系为：
\begin{equation}
  D_{xx}=D_{yy}=-\frac{1}{2}D_{zz}.
\end{equation}
将分量 $D_{zz}$ 记为D（在这种情形就简称为四极矩），我们得到势
\begin{equation}
  \varphi^{(2)}=\frac{D}{4R_0^3}(3\cos^2\theta-1)
  =\frac{D}{2R_0^3}\mathrm{P}_2(\cos\theta),
\end{equation}
式中 $\theta$ 是 $\boldsymbol{R}_0$ 与z轴之间的夹角，$\mathrm{P}_2$ 是勒让德多项式.

正如我们在上节对偶极矩所做的那样，容易证明，如果一个系统的总电荷和偶极矩都等于零，该系统的四极矩就不依赖于坐标原点的选择.

我们可以用完全相似的方法来写出展开式（41.1）中的以后各项.展开式的第l项定义了一个l阶张量（称为这个体系的 $2^l$ 极矩张量)，对于全部指标都是对称的，而当对于任何一对指标缩并时，所得结果为零；可以证明，这样的张量一共有 $2l+1$ 个独立分量.\footnote{这样的张量称为不可约张量，缩并时为零意味着不能由其分量构造出较低阶的张量.}

我们将用球谐函数理论中的著名公式
\begin{equation}
  \frac{1}{|\boldsymbol{R}_0-\boldsymbol{r}|}
  =\frac{1}{\sqrt{R_0^2+r^2-2rR_0\cos\chi}}
  =\sum_{l=0}^{\infty}\frac{r^l}{R_0^{l+1}}\mathrm{P}_l(\cos\chi)
\end{equation}
把势的展开式中的一般项表为另一种形式，式中 $\chi$ 是 $\boldsymbol{R}_0$ 和 $\boldsymbol{r}$ 之间的夹角.我们引入 $\boldsymbol{R}_0$ 和 $\boldsymbol{r}$ 分别同固定坐标轴形成的球面角 $\theta,\varPhi$ 和 $\theta',\varphi$，应用球谐函数的加法定理：
\begin{equation}
  \mathrm{P}_l(\cos\chi)=\sum_{m=-l}^{l}
  \frac{(l-|m|)!}{(l+|m|)!}
  \mathrm{P}_l^{|m|}(\cos\theta)\mathrm{P}_l^{|m|}(\cos\theta')
  \mathrm{e}^{-\mathrm{i}m(\varPhi-\varphi)},
\end{equation}
式中 $\mathrm{P}_l^m$ 是缔合勒让德多项式.

我们也引入球面函数\footnote{按照量子力学的定义.}
\begin{equation}
\begin{aligned}
  &\mathrm{Y}_{lm}(\theta,\varphi)=(-1)^m\mathrm{i}^l
  \sqrt{\frac{2l+1}{4\pi}\frac{(l-m)!}{(l+m)!}}\,
  \mathrm{P}_l^m(\cos\theta)\,\mathrm{e}^{\mathrm{i}m\varphi},\qquad m\geqslant 0,\\
  &\mathrm{Y}_{l,-|m|}(\theta,\varphi)=(-1)^{l-m}\mathrm{Y}_{l,|m|}^*.
\end{aligned}
\end{equation}
则展开式（41.9）取形式：
\begin{equation*}
  \frac{1}{|\boldsymbol{R}_0-\boldsymbol{r}|}
  =\sum_{l=0}^{\infty}\sum_{m=-l}^{l}
  \frac{r^l}{R_0^{l+1}}\frac{4\pi}{2l+1}
  \mathrm{Y}_{lm}^{*}(\theta,\varPhi)\mathrm{Y}_{lm}(\theta',\varphi).
\end{equation*}
对（40.1）中每一项进行这样的展开，我们最后得到势的展开式中第l项的如下表达式：
\begin{equation}
  \varphi^{(l)}=\frac{1}{R_0^{l+1}}\sum_{m=-l}^{l}
  \sqrt{\frac{4\pi}{2l+1}}\,Q_m^{(l)}\,\mathrm{Y}_{lm}^{*}(\theta,\varPhi),
\end{equation}
式中
\begin{equation}
  Q_m^{(l)}=\sum_a e_a r_a^l\sqrt{\frac{4\pi}{2l+1}}\,
  \mathrm{Y}_{lm}(\theta_a,\varphi_a).
\end{equation}
2l+1个量 $Q_m^{(l)}$ 的集合构成电荷系统的 $2^l$ 极矩.

以这种方式定义的量 $Q_m^{(l)}$ 与偶极矩矢量 $\boldsymbol{d}$ 的分量关系如下
\begin{equation}
  Q_0^{(1)}=\mathrm{i}d_z,\qquad
  Q_{\pm1}^{(1)}=\pm\frac{\mathrm{i}}{\sqrt{2}}(d_x\pm\mathrm{i}d_y).
\end{equation}
量 $Q_m^{(2)}$ 与张量分量 $D_{\alpha\beta}$ 的关系如下
\begin{equation}
\begin{aligned}
  &Q_0^{(2)}=-\frac{1}{2}D_{zz},\qquad
  Q_{\pm1}^{(2)}=\mp\frac{1}{\sqrt{6}}(D_{xz}\pm\mathrm{i}D_{yz}),\\
  &Q_{\pm2}^{(2)}=-\frac{1}{2\sqrt{3}}(D_{xx}-D_{yy}\pm2\mathrm{i}D_{xy}).
\end{aligned}
\end{equation}

\begin{xiti}

\noindent{\heiti 习\ 题}\quad 求一个均匀带电椭球相对于其中心的四极矩.

解：将（41.3)式中的求和换为遍历椭球体积的积分，我们有：
\begin{equation*}
  D_{xx}=\rho\iiint(2x^2-y^2-z^2)\,\mathrm{d}x\,\mathrm{d}y\,\mathrm{d}z,\quad\text{等}.
\end{equation*}
选取坐标轴沿椭球的轴，原点置于椭球中心；从对称性考虑显然可见，这些轴就是张量 $D_{\alpha\beta}$ 的主轴.借助变换
\begin{equation*}
  x=x'a,\qquad y=y'b,\qquad z=z'c,
\end{equation*}
遍历椭球
\begin{equation*}
  \frac{x^2}{a^2}+\frac{y^2}{b^2}+\frac{z^2}{c^2}=1
\end{equation*}
体积的积分就化为遍历单位球
\begin{equation*}
  x'^2+y'^2+z'^2=1
\end{equation*}
体积的积分.结果我们得到：
\begin{equation*}
\begin{aligned}
  &D_{xx}=\frac{e}{5}(2a^2-b^2-c^2),\qquad
  D_{yy}=\frac{e}{5}(2b^2-a^2-c^2),\\
  &D_{zz}=\frac{e}{5}(2c^2-a^2-b^2),
\end{aligned}
\end{equation*}
式中 $e=(4\pi/3)abc\,\rho$ 是椭球的总电荷.

\end{xiti}

\section{外场中的电荷体系}

现在我们来研究一个位于外电场中的电荷体系.我们用 $\varphi(\boldsymbol{r})$ 表示电荷 $e_a$ 所在点的外电场的势.每个电荷的势能是 $e_a\varphi(r_a)$，而这个体系的总势能则有
\begin{equation}
  U=\sum_a e_a\varphi(\boldsymbol{r}_a).
\end{equation}
我们仍然采用原点在电荷体系内一任意点上的坐标系；$\boldsymbol{r}_a$ 是电荷 $e_a$ 在这个坐标系中的径矢.

假设外场在电荷体系所在的区域中缓慢地变化，即对于该系统是准均匀的.这时，我们可以将能量U展开为 $\boldsymbol{r}_a$ 的幂级数.
\begin{equation}
  U=U^{(0)}+U^{(1)}+U^{(2)}+\cdots
\end{equation}
展开式中的第一项是
\begin{equation}
  U^{(0)}=\varphi_0\sum_a e_a,
\end{equation}
式中，$\varphi_0$ 是势在坐标原点之值.在这个近似中，体系的能量同所有的电荷都集中在一点（坐标原点）时的能量一样.

展开式中的第二项是
\begin{equation*}
  U^{(1)}=(\operatorname{grad}\varphi)_0\cdot\sum e_a\boldsymbol{r}_a.
\end{equation*}
引入原点的电场强度 $\boldsymbol{E}_0$ 和系统的偶极矩 $\boldsymbol{d}$，我们得到
\begin{equation}
  U^{(1)}=-\boldsymbol{d}\cdot\boldsymbol{E}_0.
\end{equation}
作用在准均匀外场中一个系统上的总力（精确到我们考虑的量级）是
\begin{equation*}
  \boldsymbol{F}=\boldsymbol{E}_0\sum e_a+[\operatorname{grad}(\boldsymbol{d}\cdot\boldsymbol{E})]_0.
\end{equation*}
如果总电荷为零，则第一项消失，于是有
\begin{equation}
  \boldsymbol{F}=(\boldsymbol{d}\cdot\nabla)\boldsymbol{E},
\end{equation}
就是说，力决定于场强的导数（取在原点）.作用于系统的总力矩是
\begin{equation}
  \boldsymbol{K}=\sum(\boldsymbol{r}_a\times e_a\boldsymbol{E}_0)
  =\boldsymbol{d}\times\boldsymbol{E}_0,
\end{equation}
就是说，精确到最低阶，它决定于场强本身.

假设有两个电荷体系，每一个体系的总电荷为零，而偶极矩则为 $\boldsymbol{d}_1$ 及 $\boldsymbol{d}_2$.两个体系的距离比体系本身的尺度大很多.我们来求他们的相互作用势能 $U$.为此，我们将两个体系中的一个认为处在另一个体系的场中，这时，
\begin{equation*}
  U=-\boldsymbol{d}_2\cdot\boldsymbol{E}_1,
\end{equation*}
式中，$\boldsymbol{E}_1$ 是第一个体系的场.将 $\boldsymbol{E}_1$ 的表达式（40.8）代入，我们得到
\begin{equation}
  U=\frac{(\boldsymbol{d}_1\cdot\boldsymbol{d}_2)R^2
  -3(\boldsymbol{d}_1\cdot\boldsymbol{R})(\boldsymbol{d}_2\cdot\boldsymbol{R})}{R^5},
\end{equation}
式中，R是两个体系间的距离矢量.

如果体系之一的总电荷不为零(而等于e)，则同理可得
\begin{equation}
  U=e\frac{\boldsymbol{d}\cdot\boldsymbol{R}}{R^3},
\end{equation}
式中，R是一个矢量，从偶极子指向电荷.

展开式（42.2）中的下一项是
\begin{equation*}
  U^{(2)}=\frac{1}{2}\sum ex_\alpha x_\beta
  \frac{\partial^2\varphi_0}{\partial x_\alpha\partial x_\beta}.
\end{equation*}
在这里，也像在§41中一样，我们略去了电荷编号的指标；势的二阶导数值取在原点；但是势 $\varphi$ 满足拉普拉斯方程，
\begin{equation*}
  \frac{\partial^2\varphi}{\partial x_\alpha^2}
  =\delta_{\alpha\beta}\frac{\partial^2\varphi}{\partial x_\alpha\partial x_\beta}=0.
\end{equation*}
因此，我们可以写出：
\begin{equation*}
  U^{(2)}=\frac{1}{2}\frac{\partial^2\varphi_0}{\partial x_\alpha\partial x_\beta}
  \sum e\left(x_\alpha x_\beta-\frac{1}{3}\delta_{\alpha\beta}r^2\right),
\end{equation*}
或最后有
\begin{equation}
  U^{(2)}=\frac{D_{\alpha\beta}}{6}
  \frac{\partial^2\varphi_0}{\partial x_\alpha\partial x_\beta}.
\end{equation}
级数（42.2)中的一般项可以借助上节定义的 $2^l$ 极矩 $Q_m^{(l)}$ 来表示.为此我们首先把势 $\varphi(\boldsymbol{r})$ 展开为球谐函数；这种展开的一般形式为
\begin{equation}
  \varphi(r)=\sum_{l=0}^{\infty}r^l\sum_{m=-l}^{l}a_{lm}
  \sqrt{\frac{4\pi}{2l+1}}\,\mathrm{Y}_{lm}(\theta,\varphi),
\end{equation}
式中 $r,\theta,\varphi$ 是点的球坐标，$a_{lm}$ 是常系数.构造和式（42.1）并用定义（41.13）我们得到：
\begin{equation}
  U^{(l)}=\sum_{m=-l}^{l}a_{lm}Q_m^{(l)}.
\end{equation}

\section{恒定磁场}

我们来研究一个作有限运动的电荷体系所产生的磁场.所谓有限运动，是指粒子在一切时间都在空间的有限区域内运动，而且它们的动量在一切时间都是有限的.这样的运动具有“稳定”的特征，考虑电荷所产生的磁场（对时间）的平均值 $\overline{\boldsymbol{H}}$ 是饶有兴趣的；这个平均磁场现在将仅是坐标的函数，而不是时间的函数，亦即是恒定的.

为了求出平均场 $\overline{\boldsymbol{H}}$ 的方程，我们先对麦克斯韦方程
\begin{equation*}
  \operatorname{div}\boldsymbol{H}=0,\qquad
  \operatorname{rot}\boldsymbol{H}
  =\frac{1}{c}\frac{\partial\boldsymbol{E}}{\partial t}
  +\frac{4\pi}{c}\boldsymbol{j}
\end{equation*}
取时间平均值.

这些方程中的第一个简单地给出
\begin{equation}
  \operatorname{div}\overline{\boldsymbol{H}}=0.
\end{equation}
在第二个方程中，导数 $\partial\boldsymbol{E}/\partial t$ 的平均值为零，正像在一般情况下任何一个在有限范围内变化的量的导数的平均值一样（见§34第二个脚注）.因此，第二个麦克斯韦方程化为
\begin{equation}
  \operatorname{rot}\overline{\boldsymbol{H}}=\frac{4\pi}{c}\overline{\boldsymbol{j}}.
\end{equation}
这两个方程就决定了恒定场 $\overline{\boldsymbol{H}}$

我们按照
\begin{equation*}
  \operatorname{rot}\overline{\boldsymbol{A}}=\overline{\boldsymbol{H}}
\end{equation*}
引入平均矢势 $\overline{\boldsymbol{A}}$.将这个方程代入（43.2）.可得
\begin{equation*}
  \operatorname{grad}\operatorname{div}\overline{\boldsymbol{A}}
  -\Delta\overline{\boldsymbol{A}}
  =\frac{4\pi}{c}\overline{\boldsymbol{j}}.
\end{equation*}
但是我们知道，场的矢势未被唯一地确定，因而我们可以加一个任意的附加条件.根据这点，我们这样来选择矢势 $\overline{\boldsymbol{A}}$，使
\begin{equation}
  \operatorname{div}\overline{\boldsymbol{A}}=0.
\end{equation}
这时，确定恒定磁场矢势的方程化为
\begin{equation}
  \Delta\overline{\boldsymbol{A}}=-\frac{4\pi}{c}\overline{\boldsymbol{j}}.
\end{equation}
很容易求解这个方程，因为(43.4)式与恒定电场的标势的泊松方程(36.4)完全相似，不过在那里的电荷密度 $\rho$ 被这里的电流密度 $\overline{\boldsymbol{j}}/c$ 所代替而已.与泊松方程的解（36.8）相似，我们可以直接写出
\begin{equation}
  \overline{\boldsymbol{A}}=\frac{1}{c}\int\frac{\overline{\boldsymbol{j}}}{R}\,\mathrm{d}V,
\end{equation}
式中，R 是从我们求 $\overline{\boldsymbol{A}}$ 的那一点到体积元 $\mathrm{d}V$ 的距离.

在公式（43.5）中，如果我们用 $\rho\boldsymbol{v}$ 代替 $\boldsymbol{j}$，并且记住所有的电荷都是点电荷，我们就可以从积分过渡到对电荷求和了.在这里，我们必须记着，在积分（43.5)中，R不过是一个积分变量，因此，它在求平均值的过程中不起作用. 如果我们用和
\begin{equation*}
  \sum\frac{e_a\boldsymbol{v}_a}{R_a}
\end{equation*}
代替
\begin{equation*}
  \int\frac{\boldsymbol{j}}{R}\,\mathrm{d}V,
\end{equation*}
那么，$R_a$ 就是各个粒子的径矢，这些径矢在电荷运动时是变动的.因此，我们应该写出
\begin{equation}
  \overline{\boldsymbol{A}}=\frac{1}{c}\sum\frac{\overline{e_a\boldsymbol{v}_a}}{R_a},
\end{equation}
此外我们是对求和记号下的整个式子求平均值.

知道了 $\overline{\boldsymbol{A}}$，我们就可以求出磁场，
\begin{equation*}
  \overline{\boldsymbol{H}}=\operatorname{rot}\overline{\boldsymbol{A}}
  =\operatorname{rot}\frac{1}{c}\int\frac{\overline{\boldsymbol{j}}}{R}\,\mathrm{d}V.
\end{equation*}
算符rot只关系着我们求场的那一点的坐标，因此，rot可以写在积分号内，而在微分的过程中，$j$ 可以当做常量.将熟知的公式
\begin{equation*}
  \operatorname{rot}f\boldsymbol{a}
  =f\operatorname{rot}\boldsymbol{a}
  +\operatorname{grad}f\times\boldsymbol{a}
\end{equation*}
(式中的 $f$ 和 $\boldsymbol{a}$ 是任意的标量及矢量）应用到积 $\overline{\boldsymbol{j}}\cdot\dfrac{1}{R}$，我们得到
\begin{equation*}
  \operatorname{rot}\frac{\overline{\boldsymbol{j}}}{R}
  =\operatorname{grad}\frac{1}{R}\times\overline{\boldsymbol{j}}
  =\frac{\overline{\boldsymbol{j}}\times\boldsymbol{R}}{R^3},
\end{equation*}
因此，
\begin{equation}
  \overline{\boldsymbol{H}}
  =\frac{1}{c}\int\frac{\overline{\boldsymbol{j}}\times\boldsymbol{R}}{R^3}\,\mathrm{d}V
\end{equation}
(径矢R是从dV指向我们求它的场的那一点).这就是毕奥-萨伐尔定律.

\section{磁矩}

现在我们来研究一个稳定运动着的电荷体系在与电荷体系相距很远的地方所产生的平均磁场.所谓很远，就是说与电荷体系的距离远大于电荷体系本身的尺度.

我们引入一个坐标系，让坐标原点在电荷体系内任意一点上，像在§40中所做的一样.我们仍以 $\boldsymbol{r}_a$ 代表各个电荷的径矢，而以 $\boldsymbol{R}_0$ 代表我们正在求场的那一点的径矢，那么，$\boldsymbol{R}_0-\boldsymbol{r}_a$ 就是从电荷 $e_a$ 到场点的径矢.按照（43.6）式，对于矢势我们有
\begin{equation}
  \overline{\boldsymbol{A}}
  =\frac{1}{c}\sum\overline{\frac{e_a\boldsymbol{v}_a}{|\boldsymbol{R}_0-\boldsymbol{r}_a|}}.
\end{equation}
像在§40中一样，我们将上式展开为 $\boldsymbol{r}_a$ 的幂级数.如果只要求精确到第一阶，我们有（略去指标 a）
\begin{equation*}
  \overline{\boldsymbol{A}}
  =\frac{1}{cR_0}\sum e\,\overline{\boldsymbol{v}}
  -\frac{1}{c}\sum e\,\overline{\boldsymbol{v}\left(\boldsymbol{r}\cdot\nabla\frac{1}{R_0}\right)}.
\end{equation*}
在第一项中，我们可以写出
\begin{equation*}
  \sum e\,\overline{\boldsymbol{v}}
  =\overline{\frac{\mathrm{d}}{\mathrm{d}t}\sum e\,\boldsymbol{r}}.
\end{equation*}
但是在一个有限范围内变化的量的导数的平均值（如 $\sum e\boldsymbol{r}$ 的导数的平均值）为零.因此，对于 $\overline{\boldsymbol{A}}$，余下来的式子是
\begin{equation*}
  \overline{\boldsymbol{A}}
  =-\frac{1}{c}\sum e\,\overline{\boldsymbol{v}\left(\boldsymbol{r}\cdot\nabla\frac{1}{R_0}\right)}
  =\frac{1}{cR_0^3}\sum\overline{e\,\boldsymbol{v}(\boldsymbol{r}\cdot\boldsymbol{R}_0)}.
\end{equation*}
将上式作如下变换.注意到 $\boldsymbol{v}=\dot{\boldsymbol{r}}$，我们可以写出（记着 $\boldsymbol{R}_0$ 是一个常矢量）
\begin{equation*}
  \sum e(\boldsymbol{R}_0\cdot\boldsymbol{r})\boldsymbol{v}
  =\frac{1}{2}\frac{\mathrm{d}}{\mathrm{d}t}\sum e\,\boldsymbol{r}(\boldsymbol{r}\cdot\boldsymbol{R}_0)
  +\frac{1}{2}\sum e[\boldsymbol{v}(\boldsymbol{r}\cdot\boldsymbol{R}_0)
  -\boldsymbol{r}(\boldsymbol{v}\cdot\boldsymbol{R}_0)].
\end{equation*}
当把上式代入 $\overline{\boldsymbol{A}}$ 的表达式以后，第一项（包含有对时间的导数）的平均值又为零，我们得到
\begin{equation*}
  \overline{\boldsymbol{A}}
  =\frac{1}{2cR_0^3}\sum\overline{
  e[\boldsymbol{v}(\boldsymbol{r}\cdot\boldsymbol{R}_0)
  -\boldsymbol{r}(\boldsymbol{v}\cdot\boldsymbol{R}_0)]}.
\end{equation*}
我们引入一个叫做体系的磁矩的矢量
\begin{equation}
  \mathfrak{m}=\frac{1}{2c}\sum e\,\boldsymbol{r}\times\boldsymbol{v},
\end{equation}
这样，我们就得到 $\overline{\boldsymbol{A}}$ 的表达式
\begin{equation}
  \overline{\boldsymbol{A}}
  =\frac{\overline{\mathfrak{m}}\times\boldsymbol{R}_0}{R_0^3}
  =\nabla\frac{1}{R_0}\times\overline{\mathfrak{m}}.
\end{equation}
知道了矢势，就很容易求得磁场.利用公式
\begin{equation*}
  \operatorname{rot}(\boldsymbol{a}\times\boldsymbol{b})
  =(\boldsymbol{b}\cdot\nabla)\boldsymbol{a}
  -(\boldsymbol{a}\cdot\nabla)\boldsymbol{b}
  +\boldsymbol{a}\operatorname{div}\boldsymbol{b}
  -\boldsymbol{b}\operatorname{div}\boldsymbol{a},
\end{equation*}
我们得到
\begin{equation*}
  \overline{\boldsymbol{H}}=\operatorname{rot}\overline{\boldsymbol{A}}
  =\operatorname{rot}\left(\frac{\overline{\mathfrak{m}}\times\boldsymbol{R}_0}{R_0^3}\right)
  =\overline{\mathfrak{m}}\operatorname{div}\frac{\boldsymbol{R}_0}{R_0^3}
  -(\overline{\mathfrak{m}}\cdot\nabla)\frac{\boldsymbol{R}_0}{R_0^3}.
\end{equation*}
其次，当 $R_0\neq 0$ 时，因为
\begin{equation*}
  \operatorname{div}\frac{\boldsymbol{R}_0}{R_0^3}
  =\boldsymbol{R}_0\cdot\operatorname{grad}\frac{1}{R_0^3}
  +\frac{1}{R_0^3}\operatorname{div}\boldsymbol{R}_0=0
\end{equation*}
和
\begin{equation*}
  (\overline{\mathfrak{m}}\cdot\nabla)\frac{\boldsymbol{R}_0}{R_0^3}
  =\frac{1}{R_0^3}(\overline{\mathfrak{m}}\cdot\nabla)\boldsymbol{R}_0
  +\boldsymbol{R}_0(\overline{\mathfrak{m}}\cdot\nabla)\frac{1}{R_0^3}
  =\frac{\overline{\mathfrak{m}}}{R_0^3}
  -\frac{3\boldsymbol{R}_0(\overline{\mathfrak{m}}\cdot\boldsymbol{R}_0)}{R_0^5}.
\end{equation*}
所以
\begin{equation}
  \overline{\boldsymbol{H}}
  =\frac{3\boldsymbol{n}(\overline{\mathfrak{m}}\cdot\boldsymbol{n})
  -\overline{\mathfrak{m}}}{R_0^3},
\end{equation}
式中 $\boldsymbol{n}$ 还是沿 $R_0$ 方向的单位矢量.由上式可以看出，用磁矩表示磁场的公式，像用偶极矩表示电场的公式一样（见（40.8）式）.

如果体系内的所有电荷都有相同的荷质比，那么，我们就可以写出：
\begin{equation*}
  \mathfrak{m}=\frac{1}{2c}\sum e\,\boldsymbol{r}\times\boldsymbol{v}
  =\frac{e}{2mc}\sum m\,\boldsymbol{r}\times\boldsymbol{v}.
\end{equation*}
如果所有电荷的速度 $v\ll c$，那么 $m\boldsymbol{v}$ 就是电荷的动量 $\boldsymbol{p}$，于是我们得到
\begin{equation}
  \mathfrak{m}=\frac{e}{2mc}\sum\boldsymbol{r}\times\boldsymbol{p}
  =\frac{e}{2mc}\boldsymbol{M},
\end{equation}
式中，$\boldsymbol{M}=\sum\boldsymbol{r}\times\boldsymbol{p}$ 是体系的机械角动量.因此，在现在的情况下，磁矩与角动量之比是一个常数，并且等于 $e/(2mc)$.

\begin{xiti}

\noindent{\heiti 习\ 题}\quad 求由两个电荷（速度 $v\ll c$）所组成的电荷体系的磁矩与角动量的比.

解：选择两个粒子的质心作为坐标的原点，我们便得到 $m_1\boldsymbol{r}_1+m_2\boldsymbol{r}_2=0$ 及 $\boldsymbol{p}_1=-\boldsymbol{p}_2=\boldsymbol{p}$，式中p是相对运动的动量.利用这些关系，我们得到
\begin{equation*}
  \mathfrak{m}
  =\frac{1}{2c}\left(\frac{e_1}{m_1^2}+\frac{e_2}{m_2^2}\right)
  \frac{m_1m_2}{m_1+m_2}\boldsymbol{M}.
\end{equation*}

\end{xiti}

\section{拉莫尔定理}

我们来考虑一个处于外在恒定均匀磁场中的电荷系统.

作用在该系统上的力的时间平均
\begin{equation*}
  \overline{\boldsymbol{F}}
  =\sum\frac{e}{c}\overline{\boldsymbol{v}\times\boldsymbol{H}}
  =\overline{\frac{\mathrm{d}}{\mathrm{d}t}\sum\frac{e}{c}\,\boldsymbol{r}\times\boldsymbol{H}}
\end{equation*}
是零，正如在有限范围内变化的量的时间导数的时间平均一样.力矩的平均值是
\begin{equation*}
  \overline{\boldsymbol{K}}
  =\sum\frac{e}{c}\overline{\boldsymbol{r}\times(\boldsymbol{v}\times\boldsymbol{H})}
\end{equation*}
且不等于零.通过展开矢量三重积,可以将它用系统的磁矩表示出来：
\begin{equation*}
  \boldsymbol{K}
  =\sum\frac{e}{c}\{\boldsymbol{v}(\boldsymbol{r}\cdot\boldsymbol{H})
  -\boldsymbol{H}(\boldsymbol{v}\cdot\boldsymbol{r})\}
  =\sum\frac{e}{c}\left\{\boldsymbol{v}(\boldsymbol{r}\cdot\boldsymbol{H})
  -\frac{1}{2}\boldsymbol{H}\frac{\mathrm{d}}{\mathrm{d}t}r^2\right\}.
\end{equation*}
第二项平均后为零，所以
\begin{equation*}
  \overline{\boldsymbol{K}}
  =\sum\frac{e}{c}\overline{\boldsymbol{v}(\boldsymbol{r}\cdot\boldsymbol{H})}
  =\frac{1}{2c}\sum e\{\overline{\boldsymbol{v}(\boldsymbol{r}\cdot\boldsymbol{H})}
  -\overline{\boldsymbol{r}(\boldsymbol{v}\cdot\boldsymbol{H})}\}
\end{equation*}
(最后的变换类似于推导(44.3)时用过的变换),或者最后有
\begin{equation}
  \overline{\boldsymbol{K}}=\overline{\mathfrak{m}}\times\boldsymbol{H}.
\end{equation}
请注意这个式子同电场情况的公式(42.6)类似.

处于外在恒定均匀磁场中电荷系统的拉格朗日量（同封闭系统比较）包含一个附加项
\begin{equation}
  L_H=\sum\frac{e}{c}\boldsymbol{A}\cdot\boldsymbol{v}
  =\sum\frac{e}{2c}(\boldsymbol{H}\times\boldsymbol{r})\cdot\boldsymbol{v}
  =\sum\frac{e}{2c}(\boldsymbol{r}\times\boldsymbol{v})\cdot\boldsymbol{H}
\end{equation}
(这里我们已经用了(19.4)式来表达均匀磁场中的矢势).引入系统的磁矩，我们得到：
\begin{equation}
  L_H=\overline{\mathfrak{m}}\cdot\boldsymbol{H}.
\end{equation}
我们注意到这种情形与存在电场的情形相似；在均匀电场内，一个总电荷为零的电荷体系的拉格朗日量包含有
\begin{equation*}
  L_E=\boldsymbol{d}\cdot\boldsymbol{E}
\end{equation*}
的项(d是电荷体系的偶极矩)，且在这种情形下，就等于电荷体系的势能反号（见 §42）.

我们现在考虑一个电荷体系，它在一个中心对称的电场内作有限运动(其速度 $v\ll c$)，这个中心对称场是由某一个静止电荷所产生的.

我们从静止坐标系变换到绕着通过静止电荷的轴而匀速旋转的坐标系中去.根据熟知的公式，我们得到粒子在新坐标系内的速度v与在旧坐标系内的速度 $\boldsymbol{v}'$ 关系式：
\begin{equation*}
  \boldsymbol{v}'=\boldsymbol{v}+\boldsymbol{\Omega}\times\boldsymbol{r},
\end{equation*}
式中，$\boldsymbol{r}$ 是粒子的径矢，而$\boldsymbol{\Omega}$是旋转坐标系的角速度.在静止坐标系统中，电荷体系的拉格朗日量是
\begin{equation*}
  L=\sum\frac{m\boldsymbol{v}'^2}{2}-U,
\end{equation*}
式中，$U$ 是诸电荷在外场中的势能与它们彼此之间的相互作用能之和.量U是电荷体系内各电荷与静止电荷之间的距离的函数，也是电荷体系内各电荷彼此间的距离的函数；当变换到旋转坐标时，它显然保持不变.因此，在新坐标系中，拉格朗日量将是
\begin{equation*}
  L=\sum\frac{m}{2}(\boldsymbol{v}+\boldsymbol{\Omega}\times\boldsymbol{r})^2-U.
\end{equation*}
如果所有的电荷都有同样的荷质比 $e/m$，并假定
\begin{equation}
  \boldsymbol{\Omega}=\frac{e}{2mc}\boldsymbol{H}.
\end{equation}
那么，在H足够小的情况下（当可以略去 $H^2$ 时)，拉格朗日量具有下面的形式：
\begin{equation*}
  L=\sum\frac{m\boldsymbol{v}^2}{2}
  +\frac{1}{2c}\sum e(\boldsymbol{H}\times\boldsymbol{r})\cdot\boldsymbol{v}-U.
\end{equation*}
可以看出，这个式子与存在恒定磁场时，所研究的电荷在静止坐标系中运动的拉格朗日量的表达式（见（45.2）式）一样.

因此我们得出结论：在非相对论情形下，一个电荷体系（所有电荷的荷质比 $e/m$ 都相等）在一中心对称电场和弱均匀磁场H内作有限运动，且这个电荷体系的行为与同一电荷体系在同一电场中在一个以角速度(45.4)匀速旋转的坐标系内的行为一样.这就是所谓拉莫尔定理，角速度 $\boldsymbol{\Omega}=e\boldsymbol{H}/(2mc)$ 称为拉莫尔频率.

我们可以从不同的观点来研究这同一个问题.如果磁场H足够弱，拉莫尔频率同电荷系统有限运动的频率相比很小.则我们可以考虑(在可与周期 $2\pi/\Omega$ 相比的时间内）对描述该系统的量进行平均.这些新的量将（以周期Ω)随时间缓慢变化.

我们来考虑系统角动量M平均值随时间的变化.按照力学中熟知的方程，M的导数等于作用在系统上的力矩K.因而用（45.1)，我们得到：
\begin{equation*}
  \frac{\mathrm{d}\overline{\boldsymbol{M}}}{\mathrm{d}t}
  =\overline{\boldsymbol{K}}=\overline{\mathfrak{m}}\times\boldsymbol{H}.
\end{equation*}
如果荷质比 $e/m$ 对于系统的所有粒子都相同，角动量和磁矩彼此成正比例，我们用（44.5）和（45.4）得到：
\begin{equation}
  \frac{\mathrm{d}\overline{\boldsymbol{M}}}{\mathrm{d}t}
  =-\boldsymbol{\Omega}\times\overline{\boldsymbol{M}}.
\end{equation}
这个方程表明，矢量 $\overline{\boldsymbol{M}}$ (随之磁矩$\mathfrak{m}$)以角速度$-\boldsymbol{\Omega}$绕着场的方向转动，而其绝对大小和它与这个方向所成的角保持不变.(这种运动称为拉莫尔进动.)
